\documentclass[11pt]{article}
% Shared preamble for all agentic-payments memos.
% Usage:
%   \documentclass[11pt]{article}
%   % Shared preamble for all agentic-payments memos.
% Usage:
%   \documentclass[11pt]{article}
%   % Shared preamble for all agentic-payments memos.
% Usage:
%   \documentclass[11pt]{article}
%   \input{_preamble}
%   \title{...} \author{...} \date{...}
%   \begin{document} \maketitle ... \end{document}

\usepackage[margin=1in]{geometry}
\usepackage[T1]{fontenc}
\usepackage[utf8]{inputenc}
\usepackage{lmodern}
\usepackage{microtype}
\usepackage{booktabs}
\usepackage{array}
\usepackage{longtable}
\usepackage{enumitem}
\usepackage{xcolor}
\usepackage{hyperref}
\usepackage{amsmath}
\usepackage{amssymb}
\usepackage{graphicx}
\usepackage{caption}
\usepackage{subcaption}
\usepackage{listings}

\hypersetup{
  colorlinks=true,
  linkcolor=blue!55!black,
  citecolor=blue!55!black,
  urlcolor=blue!55!black
}

\setlist[itemize]{topsep=3pt,itemsep=2pt,leftmargin=1.4em}
\setlist[enumerate]{topsep=3pt,itemsep=2pt,leftmargin=1.7em}

\definecolor{codebg}{RGB}{246,247,242}
\lstset{
  basicstyle=\ttfamily\footnotesize,
  backgroundcolor=\color{codebg},
  frame=single,
  rulecolor=\color{black!12},
  breaklines=true,
  showstringspaces=false,
  columns=fullflexible,
  keepspaces=true,
}

\newcommand{\code}[1]{\texttt{\detokenize{#1}}}
\newcommand{\risk}{\operatorname{risk}}

% Experimental result callout (used by data-driven memos).
\newenvironment{result}
  {\par\medskip\noindent\textbf{Result.}\quad\itshape}
  {\upshape\par\medskip}
\newenvironment{hypothesis}
  {\par\medskip\noindent\textbf{Hypothesis.}\quad}
  {\par\medskip}
\newenvironment{experiment}
  {\par\medskip\noindent\textbf{Experiment.}\quad}
  {\par\medskip}

%   \title{...} \author{...} \date{...}
%   \begin{document} \maketitle ... \end{document}

\usepackage[margin=1in]{geometry}
\usepackage[T1]{fontenc}
\usepackage[utf8]{inputenc}
\usepackage{lmodern}
\usepackage{microtype}
\usepackage{booktabs}
\usepackage{array}
\usepackage{longtable}
\usepackage{enumitem}
\usepackage{xcolor}
\usepackage{hyperref}
\usepackage{amsmath}
\usepackage{amssymb}
\usepackage{graphicx}
\usepackage{caption}
\usepackage{subcaption}
\usepackage{listings}

\hypersetup{
  colorlinks=true,
  linkcolor=blue!55!black,
  citecolor=blue!55!black,
  urlcolor=blue!55!black
}

\setlist[itemize]{topsep=3pt,itemsep=2pt,leftmargin=1.4em}
\setlist[enumerate]{topsep=3pt,itemsep=2pt,leftmargin=1.7em}

\definecolor{codebg}{RGB}{246,247,242}
\lstset{
  basicstyle=\ttfamily\footnotesize,
  backgroundcolor=\color{codebg},
  frame=single,
  rulecolor=\color{black!12},
  breaklines=true,
  showstringspaces=false,
  columns=fullflexible,
  keepspaces=true,
}

\newcommand{\code}[1]{\texttt{\detokenize{#1}}}
\newcommand{\risk}{\operatorname{risk}}

% Experimental result callout (used by data-driven memos).
\newenvironment{result}
  {\par\medskip\noindent\textbf{Result.}\quad\itshape}
  {\upshape\par\medskip}
\newenvironment{hypothesis}
  {\par\medskip\noindent\textbf{Hypothesis.}\quad}
  {\par\medskip}
\newenvironment{experiment}
  {\par\medskip\noindent\textbf{Experiment.}\quad}
  {\par\medskip}

%   \title{...} \author{...} \date{...}
%   \begin{document} \maketitle ... \end{document}

\usepackage[margin=1in]{geometry}
\usepackage[T1]{fontenc}
\usepackage[utf8]{inputenc}
\usepackage{lmodern}
\usepackage{microtype}
\usepackage{booktabs}
\usepackage{array}
\usepackage{longtable}
\usepackage{enumitem}
\usepackage{xcolor}
\usepackage{hyperref}
\usepackage{amsmath}
\usepackage{amssymb}
\usepackage{graphicx}
\usepackage{caption}
\usepackage{subcaption}
\usepackage{listings}

\hypersetup{
  colorlinks=true,
  linkcolor=blue!55!black,
  citecolor=blue!55!black,
  urlcolor=blue!55!black
}

\setlist[itemize]{topsep=3pt,itemsep=2pt,leftmargin=1.4em}
\setlist[enumerate]{topsep=3pt,itemsep=2pt,leftmargin=1.7em}

\definecolor{codebg}{RGB}{246,247,242}
\lstset{
  basicstyle=\ttfamily\footnotesize,
  backgroundcolor=\color{codebg},
  frame=single,
  rulecolor=\color{black!12},
  breaklines=true,
  showstringspaces=false,
  columns=fullflexible,
  keepspaces=true,
}

\newcommand{\code}[1]{\texttt{\detokenize{#1}}}
\newcommand{\risk}{\operatorname{risk}}

% Experimental result callout (used by data-driven memos).
\newenvironment{result}
  {\par\medskip\noindent\textbf{Result.}\quad\itshape}
  {\upshape\par\medskip}
\newenvironment{hypothesis}
  {\par\medskip\noindent\textbf{Hypothesis.}\quad}
  {\par\medskip}
\newenvironment{experiment}
  {\par\medskip\noindent\textbf{Experiment.}\quad}
  {\par\medskip}


\title{\textbf{Hackathon logistics}\\
\large Memo 09}
\author{Bloomsbury Tech -- Agentic Payments Hackathon Build}
\date{April 25, 2026}

\begin{document}
\maketitle

\textbf{Date:} 2026-04-25

\section{Schedule}
\begin{itemize}
  \item \textbf{Friday 7pm} --- Reception, rooftop overlooking London skyline. Networking. Selected candidates only.
  \item \textbf{Saturday 10am -- 8pm} --- Full-day build, central London. 10 hours.
\end{itemize}
\section{Theme}
AI agents that manage flow of money. Real-world problems: rent, invoicing, international payments, reconciliation.

\section{Prize}
First place = guaranteed YC interview with \textbf{Tom Blomfield} (YC Partner, Monzo founder).

\section{Judging signal}
Per Eugene's read: assessment is \textbf{founder potential via pitch}, not "best product". Pitch is \textbf{2 minutes}. Working demo is signal of velocity, not the artifact being judged.

\section{Implications}
\begin{enumerate}
  \item The pitch is the deliverable. Build supports the pitch.
  \item Spend \textbackslash{}\$\textbackslash{}geq\textbackslash{}\$1 hour rehearsing on Saturday afternoon.
  \item Optimise for \emph{one} killer slide (UMAP cluster plot with anomaly explanations) over a complete product.
  \item Frame as "Bloomsbury Tech" (existing ML lab with publications) not "we're a couple of guys with an idea". Credibility compounds with Blomfield.
  \item Skip vision; lead with the gap.
\end{enumerate}
\section{Reception strategy}
\begin{itemize}
  \item Goal: pressure-test the pitch on 2-3 fintech-savvy people. Listen for the part that makes their face register confusion or skepticism --- that's what to fix Saturday morning.
  \item Don't reveal the technical approach (UMAP, the polymarket asset reuse). Reveal the \emph{gap} (human-baseline fraud blindness).
  \item Look for potential cofounders if the cofounder bandwidth is short.
\end{itemize}

\end{document}
